\label{ch:xi_galois}
% Source document: 338

\epigraph{$\xi$ is not a free constant --- it follows necessarily from the
Galois group orders and the winding numbers of the $T^4$ torus.}{Doc.\,338}

\section{The Core Identity}

The square of the lepton mass ratio yields an exact integer
identity that both sides --- FFGFT and Galois algebra --- simultaneously
satisfy ():

\[
\boxed{
  \underbrace{\frac{(r_\mu/r_e)^2}{\xi}}_{\text{FFGFT}}
  \;=\;
  \underbrace{|\mathrm{GF}(9)^*|^2 \cdot 5^2 \cdot |\mathrm{GF}(27)|}_{\text{Galois}}
  \;=\; 8^2 \cdot 25 \cdot 27 \;=\; 43200
}
\]

The three Galois factors have concrete algebraic significance:
\begin{itemize}
  \item $8 = |\mathrm{GF}(9)^*| = 3^2 - 1$: multiplicative group of the 2nd generation
  \item $27 = |\mathrm{GF}(27)| = 3^3$: field of the 3rd generation (all elements)
  \item $5 = $ smallest prime divisor of $|\mathrm{GF}(81)^*| = 80$:
    first new prime in $\mathrm{GF}(3^4)^*$
\end{itemize}

\section{Derivation of \texorpdfstring{$\xi$}{xi} from the Identity}

Since both sides are the same number, $\xi$ follows algebraically:
\[
  \xi = \frac{(r_\mu/r_e)^2}{43200}
  = \frac{(12/5)^2}{43200}
  = \frac{144/25}{43200}
  = \frac{1}{7500}
  = \frac{4}{30000}
\]

The winding numbers $(r_e = 4/3,\; r_\mu = 16/5)$ in turn follow
from the Galois group structures of $\mathrm{GF}(3^n)$:

\begin{center}
\begingroup\small
\begin{tabular}{@{}p{1.2cm}p{0.8cm}p{2.8cm}p{2.8cm}@{}}
\toprule
Quantity & Value & Galois Origin & Remark \\
\midrule
$n_{\theta,1}$ & 3 & $\mathrm{char}(\mathrm{GF}(3^n))$ & axiomatic \\
$n_{\phi,1}$   & 2 & $|\mathrm{GF}(3)^*|$ & mult.\ group of the prime field \\
$n_{\theta,2}$ & 5 & recursion $3+2$ & smallest prime divisor $|\mathrm{GF}(81)^*|=80$ \\
$n_{\phi,2}$   & 4 & $\varphi(5)$ & Euler $\varphi$ of the prime divisor \\
$n_{\theta,3}$ & 9 & recursion $5+4$ & simultaneously $\mathrm{char}^2$ \\
$n_{\phi,3}$   & 5 & smallest prime divisor & \\
\bottomrule
\end{tabular}
\endgroup
\end{center}

Recursion law: $n_{\theta,g} = n_{\theta,g-1} + n_{\phi,g-1}$.

\section{Connection to \texorpdfstring{$\alpha$}{alpha}}

The geometric correction factor $K_\mathrm{frak}$ cancels exactly in
mass ratios. It connects $\xi$ with the
fine structure constant (Doc.\,070):
\[
  \alpha = \frac{\xi \cdot E_0^2}{K_\mathrm{frak}}, \quad
  E_0 = \sqrt{m_e m_\mu} \approx 7{,}35\;\text{MeV}
\]
\[
  \Rightarrow\quad \alpha^{-1} = 137{,}06 \quad (\text{Dev.\,0{,}02\,\%})
\]

Thus $\xi$ is anchored twice:
\begin{enumerate}
  \item \textbf{Algebraically} from the Galois group orders $\mathrm{GF}(3^n)$
    and the torus winding numbers --- purely structural, no free parameter.
  \item \textbf{Empirically} through $\alpha$: A single measured value
    ($\alpha \approx 1/137{,}036$) anchors $\xi$ in physical reality.
\end{enumerate}

\section{Comparison: FFGFT vs.\,Galois vs.\,Experiment}

\begin{center}
\begingroup\small
\begin{tabular}{@{}p{3.5cm}p{1.5cm}p{1.8cm}p{1.8cm}@{}}
\toprule
Method & $(m_\mu/m_e)^2$ & $m_\mu/m_e$ & Dev.\,exp. \\
\midrule
Experiment (PDG) & 42\,753 & 206{,}768 & --- \\
FFGFT bare & 43\,200 & 207{,}846 & $+0{,}52$\,\% \\
Galois $8^2\cdot25\cdot27$ & 43\,200 & 207{,}846 & $+0{,}52$\,\% \\
\bottomrule
\end{tabular}
\endgroup
\end{center}

FFGFT and Galois give exactly the same bare value. The 0{,}52\,\%
deviation from experiment is the same fractal-recursive
higher-order correction in both frameworks --- it is already fully contained in the
PDG measured values.

\section{Significance for the Hyperbit Discussion}

Matzke's hyperbit begins with $e_i^2 = +1$ --- a structural
stipulation whose justification remains open. The FFGFT counterpart is not
$\xi$ (which is algebraically derived), but the choice of the field
$\mathrm{GF}(3^n)$ as the fundamental structure. This choice is motivated by the
three-generation structure of fermions, but is not yet a
complete derivation.

The decisive difference: In the FFGFT, the step from the
algebra ($\mathrm{GF}(3^n)$, winding numbers) to the numerical value $\xi$ is closed.
In Matzke's framework, the step from $e_i^2 = +1$ to the bit mass
($m_{\text{bit}} = m_P \ln 2 / (2\pi)$) remains a borrowed bridge via Landauer
at Planck temperature --- no algebraic derivation.